%==============================================================================
% 引言
%==============================================================================
\section{引言}\label{sec:intro}

%------------------------------------------------------------------------------
% 研究背景
%------------------------------------------------------------------------------
\subsection{研究背景}

飞行器轨迹跟踪控制是航空领域的核心问题之一，广泛应用于自主飞行、编队协同、空中加油等任务场景\cite{Qin2003,Mayne2014}。在实际飞行中，飞行器动力学具有强非线性、多变量耦合以及幅值约束等特征，例如推力受发动机性能限制、航向角速率受结构强度约束。这些因素使得传统线性控制方法难以在复杂机动条件下兼顾跟踪精度与约束满足。

模型预测控制（Model Predictive Control, MPC）作为一种基于优化的控制策略，能够显式处理输入与状态约束，并通过滚动优化实现对未来行为的预测\cite{Garcia1989,Rawlings2009}. 自20世纪70年代提出以来，MPC已在过程工业\cite{Qin2003}、汽车自动驾驶\cite{Falcone2007}以及航空航天等领域得到广泛应用。Mayne\cite{Mayne2014}对MPC的最新发展与未来前景进行了系统综述，指出MPC在处理复杂约束系统方面具有独特优势。

%------------------------------------------------------------------------------
% 文献综述
%------------------------------------------------------------------------------
\subsection{相关工作}

在飞行器轨迹跟踪领域，研究者们提出了多种基于MPC的控制方案。根据对非线性动力学的处理方式，主要可分为以下三类：

\textbf{（1）非线性模型预测控制（NMPC）。} NMPC直接利用非线性动力学模型进行预测优化，理论上可获得最优控制性能。Chen和Allg\"ower\cite{Chen1998}提出了准无穷时域NMPC方法，通过引入终端约束和终端代价保证稳定性。然而，NMPC每步需求解非线性规划问题，计算量大，难以满足实时性要求。对于飞行器等快动态系统，NMPC的实时实现仍面临挑战。

\textbf{（2）线性时变模型预测控制（LTV-MPC）。} LTV-MPC通过在参考轨迹附近进行Jacobian线性化，将非线性预测问题转化为一系列二次规划（QP）问题\cite{Borrelli2017}. 该方法兼顾了非线性动力学的局部精度与QP的高效求解，在实时性与精度之间取得了良好平衡。Falcone等\cite{Falcone2007}将LTV-MPC应用于车辆主动转向控制，验证了其实时可行性。

\textbf{（3）编队飞行控制。} 多飞行器编队是单机跟踪的自然扩展。Keviczky等\cite{Keviczky2008}提出了一种分散式滚动时域控制方法，用于自主飞行器编队的协同控制。Richards和How\cite{Richards2004}研究了基于MPC的多UAV协同方法。然而，现有编队控制研究多关注防撞与协调策略，对编队内个体轨迹跟踪精度与实时性的综合评估仍有不足。

%------------------------------------------------------------------------------
% 本文贡献
%------------------------------------------------------------------------------
\subsection{本文贡献}

针对上述问题，本文提出一种基于LTV-MPC的飞行器轨迹跟踪控制方法，并扩展至多机编队场景。主要贡献如下：

\begin{enumerate}[label=（\arabic*）,leftmargin=*]
    \item \textbf{实时LTV-MPC求解框架。} 基于三自由度质点模型，通过Jacobian线性化与QP重构，将非线性预测问题转化为标准QP形式并利用quadprog高效求解。相比基于fmincon的非线性求解器，求解速度提升约85倍（从138\,ms降至1.42\,ms），满足实时控制要求。
    
    \item \textbf{参考轨迹生成方法。} 设计了基于弧长重参数化与转角平滑的轨迹生成算法，统一处理航向角unwrap与重力补偿问题，并对减速板引入的负推力进行建模，使生成的参考轨迹在动力学上可行。
    
    \item \textbf{编队扩展与验证。} 将单机LTV-MPC扩展为Leader-Follower编队架构，在标称工况下对3机编队进行了仿真验证，分析了编队跟踪误差分布。
\end{enumerate}

本文其余部分组织如下：第\ref{sec:model}节建立三自由度动力学模型；第\ref{sec:mpc}节详述LTV-MPC控制器设计；第\ref{sec:trajectory}节介绍参考轨迹生成方法；第\ref{sec:experiments}节给出仿真实验与结果分析；第\ref{sec:conclusion}节总结全文并展望未来工作。
